\scanpage{8}
as recalled above\deleted{\unsure{x}}, the morphism $\underline{\operatorname{Pic}}_{Z_n\mathord{\tikz[baseline=-.25ex,x=1em,y=1em]{\draw[line width=.7pt] (0,-.05)--(.32,.72);\fill (.13,.45) ellipse[x radius=.06em,y radius=.08em];\fill (.14,.18) ellipse[x radius=.06em,y radius=.08em];}}k}\handwritten{\longrightarrow}\underline{\operatorname{Pic}}_{X_o/k}$ is affine, and hence $\operatorname{NS}(Z_n)\handwritten{\longrightarrow}\operatorname{NS}(X_o)$ has finite kernel.) I expect that this is always the case, \deleted{and it}\tikz[baseline=(s.base)]{\node[inner sep=0pt](s){xx};} In case $A$ is normal of dimension $2$, this follows from Du-Val - Mumford's negative definiteness. It would he nice to have a proof in general( perhaps through a negative definiteness of an intersection matric $X^i.C^j$, where for every irreducible component $X^i$ of $X_o$, $C^i$ denotes \tikz[baseline=(s.base)]{\node[inner sep=0pt](s){a curve};\node[anchor=south,inner sep=0pt] at ([xshift=-.6em,yshift=-1pt]s.north){\inserted{suitable}};} on $X^i$, for instance defined in terms of a suitable \deleted{ample} relatively ample sheaf on $X$ by taking intersections of corresponding hypersurface sections on the $X^i$'\deleted{\unsure{s}} ?).

In case (12) is indeed injective, we have no trouble defining
\[
\begin{aligned}
(13)\qquad\qquad & Q=P/\underline{D}
\end{aligned}
\]
as a \uline{group-scheme \rlap{r}l\rlap{e}ocally of finite type over} $k$. We then have
\[
\begin{aligned}
(14)\qquad\qquad & P^o\handwritten{\xrightarrow{\sim}}Q^o\ ,
\end{aligned}
\]
an isomorphism for the connected components, and
\[
\begin{aligned}
(15)\qquad\qquad & Q/Q^o=(P/P^o)/\underline{D}\handwritten{\simeq}\underline{\operatorname{NS}}_{Z_n/k}/\underline{D}\ ,
\end{aligned}
\]
where $n$ is chosen large enough. Of course (14) and (11) yield an i\rlap{x}somorphism
\[
\begin{aligned}
(16)\qquad\qquad & \underline{\operatorname{Lie}}(Q)\handwritten{\simeq}H^1(X,\underline{O}_X)\ .
\end{aligned}
\]
Let us remark by the way that we have an \uline{injective} restriction morphism
\[
\begin{aligned}
(17)\qquad\qquad & H^1(X,\underline{O}_X)\handwritten{\hookrightarrow}H^1(U,\underline{O}_U)=H_s^2(S)\qquad\qquad\mathord{\tikz[baseline=-.3ex,x=1em,y=1em]{\fill (.1,0) circle (.35pt);\fill (.3,.3) circle (.35pt);}}
\end{aligned}
\]
as one sees by remarking that, with the above notations for $\underline{J}$, $\underline{L}$, \inserted{\handwritten{the sheaf}} \deleted{\unsure{of}} $i_{\handwritten{*}}(\underline{O}_U)/\underline{O}_X$ $(i:U\handwritten{\longrightarrow}X$ the inclusion) has a composition series whose quotients are the $\underline{L}^{-n}$, $n\handwritten{\geq}1$, whose $H^o$ is zero, and the map (17) is just the canonical map $H^1(X,\underline{O}_X)\mathrel{\tikz[baseline=(s.base)]{\node[inner sep=0pt](s){$\unsure{\longrightarrow|\hookrightarrow}$};\path[fill=gray,opacity=.45] ([xshift=2pt,yshift=2pt]s.west)--([xshift=3pt,yshift=12pt]s.west)--([xshift=6pt,yshift=8pt]s.west)--([xshift=-1pt,yshift=2pt]s.east)--cycle;\draw[gray,line width=.6pt] ([xshift=4pt,yshift=11pt]s.west) to[out=250,in=100] ([xshift=3pt,yshift=5pt]s.west);}}H^1(X,i_{\handwritten{*}}(\underline{O}_U))$. Th\rlap{i}us, the Lie algebra of the local Picard scheme is identified with a subspace of $H^1(U,\underline{O}_U)$. (This suggests that if $A$ is Cohen-Macaulay of dim.$\handwritten{\geq 3}$, $Q$ is discrete, hence $\operatorname{Pic}(U)$ countable etc ).
% archive sheet number: 7
